\chapter{The Capstone: Decisions and Trade-offs}
\label{ch:capstone}

\begin{quote}\itshape
The code is the easy half. The section explaining why it is shaped this way is
the half that gets read.
\end{quote}

\noindent
The mini-project is finished by the time you reach this chapter; it was built a
stage at a time across the preceding sixteen. This chapter deliberately
introduces no new API. It walks the completed system as a set of decisions --- what
was chosen, what was rejected, what it cost, and what would have to change for
the rejected option to win.

\chapterstub{
\textbf{Argument.} A system is best explained by its rejected alternatives. The
chapter is therefore organised as a decision log rather than a tour of the
codebase.

\textbf{Sections.} (17.1) The finished system end to end, at a level someone
could rebuild from. (17.2) The decision log, each entry stating the choice, the
alternative, the reason and the cost: why single-agent rather than a supervisor;
why \code{durability="sync"} on one path and \code{"async"} on the rest; why an
idempotency key rather than a distributed transaction; why SSE rather than
WebSockets; why Postgres rather than Redis; why retrieval as a tool rather than
a fixed step. (17.3) What the numbers came out at: cost per conversation, p95
latency, evaluation scores, tokens per turn --- collected from the measurements
made throughout the book rather than asserted here. (17.4) What broke while
building it, honestly. (17.5) What is deliberately missing and what it would take
to add. (17.6) Turning the project into something legible to a reviewer or an
interviewer --- what a decisions section should contain and what makes one
worthless.

\textbf{Listings.} No new API surface. Only assembled excerpts already introduced
elsewhere, shown in their final form.

\textbf{Diagrams.} \code{capstone-architecture} --- the finished system.
\code{capstone-graph} --- the actual compiled graph, rendered from the running
code rather than drawn by hand. \code{capstone-decision-log} --- the decisions as
a table-diagram with their alternatives.

\textbf{Mini-project.} Final stage: the complete service, plus the decisions
document that is the point of the whole exercise.

\textbf{Source topics covered.} Part~IX 9.2; and the collected measurements from
Chapters~\ref{ch:asyncio}, \ref{ch:tools}, \ref{ch:multi-agent},
\ref{ch:context} and \ref{ch:serving}.
}
